% =============================================================================
% 4.4 本章小结
% 草稿版本：v0.1 (2026-05-07)
% 起草动机：按导师意见 10，第 4 章必须包含独立的"本章小结"小节，承担
%           4.1 双通路 RAG / 4.2 代码生成与受限沙箱 / 4.3 自适应报告
%           生成三组件的统一收口，并把模块层与集成层的核心数据汇成
%           一段集中证据，引出第 5 章实验验证。原 4_3_报告生成.tex
%           末段的章节过渡块迁出到本文件并扩写为独立小节。
% 字数：约 750 字
% 风格规则：v0.2 风格（无 \textbf / \textit / 引例括号 / 双引号 / 破折号）
% 引用：暂无新增引用（核心数据来自 4.1 / 4.2 / 4.3 节及第 5 章评测）
% 图表：无新增（核心数据已在前 3 节图表中给出，本节仅以叙述方式收口）
% 依赖宏包：booktabs / tabularx（主稿已加载）
% =============================================================================

\subsection{本章小结}
\label{sec:chap4-summary}

本章针对开题报告中提出的异构数据语义化与数值理解幻觉两个研究问题，
构建了由双通路 RAG、代码生成与受限沙箱、自适应报告生成三个组件协同
的完整气象数据统计分析方法。\ref{sec:semantic-bridge} 节解决了从数值
到带引用源的语义标签这一离散映射任务，\ref{sec:code-execution} 节
解决了从原始数据到统计结果这一连续计算任务，\ref{sec:report-generation}
节解决了把事实数据、语义解读、统计结果三类异质输入编排为场景化决策
报告这一信息整合任务。三组件在数据流向上首尾相接，分别对应大语言
模型在精确归属、精确计算、专业表达三类能力上的内在边界，并以确定性
组件外置的工程范式将这些边界转化为可量化、可控的工程问题。

\ref{sec:semantic-bridge} 节通过双通路 RAG 架构把概念检索与数值
归属分离。通路 A 在 60 条混合检索评测集上以 $\alpha = 0.9$ 取得
Top-1 命中率 85.0\% 与 MRR 0.908，并量化了最优权重随知识库体量
漂移的工程结论；通路 B 在 71 条覆盖 9 类要素的语义桥接评测集上把
\texttt{grade\_\bk{}id} 准确率与 \texttt{citation} 真实性同时提升到
100\%，相对让大语言模型直接处理裸数据的 \texttt{llm\_\bk{}baseline}
基线在 \texttt{grade\_\bk{}id} 严格命中率上形成 94.4 个百分点的差距。
桥接前后模型上下文对比（\reffig{fig:bridge-context-compare}）从可视
角度直接呈现了这一组件的工程价值。

\ref{sec:code-execution} 节通过结构化数据契约约束的代码生成与白名单
受限沙箱，把统计计算从大语言模型的自由发挥改造为确定性 Python 函数
执行。在 18 条覆盖最大值、平均值、阈值次数等 6 类计算任务的评测中，
\texttt{llm\_\bk{}with\_\bk{}code} 档数值准确率达到 88.9\%，比让大语言
模型直接给出数值答案的 \texttt{llm\_\bk{}direct} 档高出 22.2 个百分点。
失败用例集中暴露了大语言模型在伴随状态维护与离散计数任务上的
内在边界，这一边界与 \ref{sec:semantic-bridge} 节关于数值阈值排序
的边界共同构成本文确定性组件外置范式的两个量化锚点。

\ref{sec:report-generation} 节通过场景驱动的 System Prompt 与三条
不变性约束（事实优先、出处必引、场景对齐），把多源信息整合为可读、
可溯、可决策的自然语言报告。在 30 条覆盖 6 类场景的端到端评测中，
\texttt{full} 档相对去 RAG 的 \texttt{ablated} 档把任务通过率从
73.3\% 提升到 86.7\%，把权威标准引用率从 0\% 提升到 26.7\%，把
\texttt{decision\_\bk{}support} 类用例通过率从 50\% 提升到 100\%。
四维度 LLM-as-judge 评分进一步显示，事实性、完整性、流畅性三档接近
持平，差距全部集中在专业性维度的 $+0.37$ 分差，印证了双通路 RAG
的价值在专业表达层面的精准定位。

三组件的实证数据共同支撑本章方法论的有效性。下一章将在统一实验环境
下完成 5 套独立评测集共 214 条用例的系统化验证，按模块层与集成层
两个粒度系统报告全部评测设计、量化数据与失败模式分析。
